\documentclass[12pt]{article}
\usepackage[margin=1in]{geometry}
\usepackage{amsmath}
\usepackage{natbib}
\usepackage{booktabs}
\usepackage{hyperref}

\title{Sovereign Immunity and the Enforcement of State Debts:\\
From \emph{Chisholm} to the Defaults of the 1840s}
\author{George J. Hall and Thomas J. Sargent\thanks{Background memorandum
prepared with Claude for the authors' fiscal-history project. Claims marked
\texttt{[PRIMARY]} in the accompanying \texttt{eleventh\_amendment\_state\_debts.bib}
were checked against primary texts (the case reports, the engrossed amendment,
and the NBER/AER papers with extracted tables). Figures and quotations that could
not be independently confirmed are collected in Section~\ref{sec:caveats}. This
memorandum is a companion to our notes on the Revolutionary British debts and on
the constitutional debt reforms of the 1840s.}}
\date{\today}

\begin{document}

\maketitle

\begin{abstract}
In \emph{Chisholm v. Georgia} (1793)---a suit to collect a Revolutionary War
supply debt owed by Georgia---the Supreme Court held that a state could be sued in
federal court by a citizen of another state. The states, carrying large unpaid war
debts and fearing a flood of creditor suits, reacted with extraordinary speed: a
proposed amendment reached the Senate the day after the decision, Congress
proposed the Eleventh Amendment in 1794, and it was ratified in 1795. The
Amendment restored state sovereign immunity, barring federal suits against a state
by citizens of another state or by foreign subjects; \emph{Hans v. Louisiana}
(1890)---itself a bond-repudiation case---later extended the bar even to a state's
own citizens. The consequence for public finance was decisive and durable: a
holder of a state's bonds, foreign or domestic, had \emph{no} judicial remedy to
compel payment. When nine American governments defaulted in 1841--1843 and five
repudiated all or part of their debts, European creditors---who held perhaps half
of the roughly \$198 million of state debt---discovered exactly this. With no
court to turn to, they could enforce repayment only through reputation and the
threat of exclusion from capital markets: the Rothschilds refused the federal
government ``a dollar,'' Baring would float no American paper while any state
remained in default, and Sydney Smith lampooned Pennsylvania in print. That the
paying states resumed while the repudiators were shut out is the clearest evidence
that American state debt was disciplined by reputation, not by law---a nineteenth-
century laboratory for the modern reputational theory of sovereign debt.
\end{abstract}

\section{Introduction: What Enforces a State's Debt?}

A recurring question in our project is what makes a government's promise to pay
credible when the creditor cannot compel performance. For debts of the
\emph{United States}, the answer that emerged after 1790 combined a funded revenue,
a reputation painstakingly built, and---as our companion note on the Revolutionary
British debts describes---the enforcement of treaty obligations in the federal
courts. For debts of the \emph{states}, the answer was different, because a
constitutional amendment deliberately removed the courts from the picture. This
memorandum traces how that came about, beginning with a lawsuit over a
Revolutionary supply debt, and how it shaped the way foreign creditors regarded
American state bonds before and after the defaults of the early 1840s.

\section{Chisholm v. Georgia (1793)}
\label{sec:chisholm}

The precipitating case grew out of the Revolution. In October 1777 Georgia's
commissioners bought supplies for its troops from Robert Farquhar, a Charleston
merchant; Georgia took the goods but never paid. After Farquhar's death, his
executor, Alexander Chisholm---a citizen of South Carolina---sought payment,
was rebuffed by the Georgia legislature, and eventually sued the State directly in
the Supreme Court under its original jurisdiction. Georgia refused even to appear,
asserting that a sovereign state ``cannot be drawn or compelled'' to answer
\citep{chisholm1793, mathis1967chisholm}.

Decided 18 February 1793, \emph{Chisholm} held, four to one, that Article~III
extended the federal judicial power to a suit against a state by a citizen of
another state---so that a state could be sued without its consent. The Justices
wrote seriatim: Chief Justice Jay and Justices Wilson, Blair, and Cushing in the
majority, with Wilson advancing a strong theory of national popular sovereignty
(``do the people of the United States form a Nation?''); Justice Iredell dissented
alone, on the ground that no statute authorized compelling a state to answer
\citep{chisholm1793}. The decision alarmed the states at once. They carried large
unpaid Revolutionary War debts, had confiscated Loyalist property, and had issued
and in places repudiated paper money; \emph{Chisholm} threatened to open the
federal courts to a flood of suits by creditors---out-of-state and British and
Loyalist alike---against the states \citep{orth1987judicial, jacobs1972eleventh}.
(The Farquhar claim itself was never enforced against Georgia; the family pursued
it for decades, and the Georgia legislature finally paid it by a relief act in
1847---some seventy years after the goods were delivered.)

\section{The Eleventh Amendment and State Sovereign Immunity}
\label{sec:amendment}

The reaction was among the swiftest in constitutional history. A resolution
proposing an amendment reached the Senate on 19~February 1793, the day after the
decision. Congress formally proposed the Eleventh Amendment on 4~March 1794,
having passed it in the Senate 23--2 (14~January 1794) and the House 81--9
(4~March 1794); it was ratified on 7~February 1795, though its adoption was not
proclaimed until President Adams's message of 8~January 1798
\citep{amend11_1795}. The text is spare:
\begin{quote}
The Judicial power of the United States shall not be construed to extend to any
suit in law or equity, commenced or prosecuted against one of the United States by
Citizens of another State, or by Citizens or Subjects of any Foreign State.
\end{quote}

The Amendment overruled \emph{Chisholm} and restored state sovereign immunity in
the federal courts, removing federal jurisdiction over two classes of suit against
a state: those brought by citizens of \emph{another} state, and those brought by
\emph{foreign} citizens or subjects. The practical upshot for public finance is
the load-bearing fact of this memorandum: a private creditor from another
state---\emph{or a foreign creditor}---could not sue a state in federal court to
enforce a debt, including a defaulted state bond.

Two later decisions closed the remaining gaps, and it is telling that both arose
from state-debt repudiation. In \emph{Hans v. Louisiana} (1890), a Louisiana
citizen sued his own state for unpaid interest on its ``consolidated bonds'' after
Louisiana amended its constitution to stop paying; the Court held that a state may
not be sued in federal court even by its own citizens without consent, extending
immunity beyond the Amendment's literal text \citep{hans1890, orth1987judicial}.
And in \emph{Principality of Monaco v. Mississippi} (1934), arising from repudiated
Mississippi bonds that had come into Monaco's hands, the Court held that a foreign
\emph{state}, too, cannot sue a U.S. state without its consent \citep{monaco1934}.
As John Orth shows, the modern meaning of the Eleventh Amendment was largely worked
out in precisely these debt cases \citep{orth1987judicial}.

\section{The Consequence: Debt Enforceable by Reputation, Not Law}
\label{sec:consequence}

The combined effect of the Eleventh Amendment and \emph{Hans} was that holders of
state bonds---whether from another state, a foreign country, or the issuing state
itself---had essentially no judicial remedy to compel a state to pay. As
\citet{english1996understanding} puts it, the American state debts of the period
are ``properly viewed as sovereign debts both because the U.S. Constitution
precludes suits against the states to enforce payment and because most of the
debts were held by residents of other states or countries.'' Enforcement rested,
therefore, on a state's \emph{willingness} to pay---driven by reputation and the
value of continued access to capital markets---rather than on any court's power to
compel it \citep{roberts2010ungovernable, jenks1938migration}. This is what makes
the 1840s a natural experiment in the reputational theory of sovereign debt.

\section{The Borrowing Boom and the Defaults of 1841--1843}
\label{sec:defaults}

During the 1820s and 1830s the states borrowed heavily---largely in European
capital markets---to finance canals, railroads, and state-chartered banks. By 1841
state debts outstanding totaled about \$198 million, more than half of it
authorized since the Panic of 1837, and a large share held abroad
\citep{wallis2004sovereign, kimwallis2005market}. The paper was placed through the
great transatlantic houses---Baring Brothers (with the U.S. agents Thomas Wren Ward
and Prime, Ward \& King), N.\,M. Rothschild \& Sons (through August Belmont), and
Hope \& Co.\ of Amsterdam---and, for the southwestern bank bonds, through Nicholas
Biddle's Bank of the United States of Pennsylvania, which contracted in 1838 to
market the Mississippi Union Bank bonds in Europe \citep{wallis2004sovereign}.

The depression that followed the Panic of 1837 broke this system. Between 1841 and
1843 eight states and the Florida Territory---nine governments responsible for
roughly two-thirds of all American government debt in private hands---defaulted
(Table~\ref{tab:defaults}). Five of the nine repudiated all or part of their
debts, and, tellingly, in every case the repudiation was tied to borrowing for
\emph{banks}: Mississippi and Florida repudiated their bank bonds outright,
Arkansas repudiated the Real Estate Bank (Holford) bonds, Michigan repudiated the
\$2.3 million of its ``\$5 million loan'' for which it claimed it had received no
payment, and Louisiana repudiated \emph{de facto} by never paying interest on some
\$21 million of bonds issued for three land banks \citep{wallis2004sovereign,
english1996understanding}. The other four defaulters---Indiana, Illinois, Maryland,
and Pennsylvania---suspended interest temporarily and eventually resumed. Crucially,
default was not fiscally inevitable: Mississippi's ordinary tax revenues alone would
have serviced the Union Bank bonds, so its repudiation was a choice
\citep{wallis2004sovereign}. (Our companion note treats the constitutional debt
limits that this crisis provoked.)

\begin{table}[htbp]
\centering
\caption{The nine defaulting governments, 1841--1843}
\label{tab:defaults}
\small
\begin{tabular}{l r l l}
\toprule
Government & Debt, 1841 & Default (date) & Outcome \\
\midrule
Indiana            & \$12{,}751{,}000 & Jan.\ 1841 & Resumed \\
Florida (Territory) & \$4{,}000{,}000  & Jan.\ 1841 & \textbf{Repudiated} (bank bonds) \\
Mississippi        & \$7{,}000{,}000  & Mar.\ 1841 & \textbf{Repudiated} (Union \& Planters) \\
Arkansas           & \$2{,}676{,}000  & Jul.\ 1841 & \textbf{Repudiated} in part (Holford) \\
Michigan           & \$5{,}611{,}000  & Jul.\ 1841 & \textbf{Repudiated} \$2.34M \\
Illinois           & \$13{,}527{,}292 & Jan.\ 1842 & Resumed \\
Maryland           & \$15{,}214{,}761 & Jan.\ 1842 & Resumed \\
Pennsylvania       & \$33{,}301{,}013 & Aug.\ 1842 & Resumed \\
Louisiana          & \$23{,}985{,}000 & Feb.\ 1843 & \textbf{Repudiated} \emph{de facto} \\
\bottomrule
\end{tabular}

\smallskip
\footnotesize Source: \citet{wallis2004sovereign}, Tables 1--2 (default dates from
\citealp{english1996understanding}). Ohio, often mistakenly listed, did \emph{not}
default. Total state debt outstanding in 1841 was about \$198 million.
\end{table}

\section{Foreign Creditors Before and After the Defaults}
\label{sec:creditors}

\paragraph{Before.} In the 1820s and 1830s European investors lent to American
states freely, drawn by yields of six or seven percent against about three percent
on their own governments' bonds, and reassured by the young republic's hard-won
credit. Some lent in the belief or hope that the federal government stood behind
the states, recalling Hamilton's assumption of the state Revolutionary debts in
1790 \citep{wallis2004sovereign, roberts2010ungovernable}.

\paragraph{After.} The defaults and repudiations shattered that confidence, and
because no creditor could sue, the reaction played out entirely in the market and
the press. In 1842 the \emph{federal} government---though it had extinguished its
own debt in 1835---could not place a loan in London or Amsterdam, tainted by the
states' conduct; James de Rothschild reportedly told the American agent, ``You may
tell your government that you have seen the man who is at the head of the finances
of Europe, and that he has told you that they cannot borrow a dollar, not a
dollar'' \citep[39--40]{sexton2005debtor}. Baring instructed its agent that ``no
new loan shall be introduced here while there is any one of the states as a
defaulter'' \citep[39]{sexton2005debtor}. Sydney Smith, a canon of St.\ Paul's who
held repudiated Pennsylvania bonds, published his \emph{Letters on American Debts}
(1843), confessing ``a disposition to seize and divide'' any Pennsylvanian he met
at a London dinner \citep[59]{mcgrane1935foreign} and petitioning Congress against
``men who prefer any load of infamy, however great, to any pressure of taxation,
however light'' \citep[9]{sydneysmith1843letters}. Wordsworth wrote protest
verses; the \emph{Times} declared Europeans would not lend to America again for
fifty years; and the repudiating states answered in kind, Mississippi's governor
denouncing ``Baron Rothschild, with the blood of Judas and Shylock in his veins''
\citep{mcgrane1935foreign}.

Foreign creditors and their bankers lobbied instead for the one remedy the
courts could not supply: \emph{federal assumption} of the state debts, expressly
on the Hamiltonian precedent of 1790. Baring circulated a memorandum arguing that
``a national pledge would undoubtedly collect capital together from all parts of
Europe,'' and Hope \& Co.\ pressed the same case; a Whig-sponsored assumption
proposal was debated in Congress in 1843. Democrats denounced it as ``a foreign
plot'' to fasten a national debt on a country that had just retired one, and it
went nowhere \citep{wallis2004sovereign, mcgrane1935foreign, roberts2012americas}.

\section{Reputation, Not Law: The Evidence and the Legacy}
\label{sec:reputation}

That the crisis was resolved without courts, and largely without federal
intervention, is what makes it evidence for the reputational view. At first the
market did not discriminate: the credit of \emph{all} the states, and even of the
never-defaulting federal government, was impaired. But over the following years
lenders sharply distinguished conduct. \citet{english1996understanding} shows that
the states which resumed payment were able to borrow again in the run-up to the
Civil War, while the repudiators---Mississippi, Florida, Arkansas---were, for the
most part, shut out. Since no legal sanction was available, this differential
repayment is direct evidence that states paid in order to preserve their
reputations and their access to international capital, exactly as reputational
models of sovereign debt predict. That most heavily indebted states chose to raise
taxes and pay, rather than repudiate, points the same way
\citep{english1996understanding, wallis2004sovereign}.

The episode's institutional legacy reinforced the lesson. Between 1842 and 1852
eleven states wrote new constitutions embedding procedural debt restrictions---
requiring that taxes be levied when bonds were issued, and that voters approve them
in referenda---and restricting state investment in private corporations
\citep{wallis2005constitutions}. Having learned that neither the courts nor the
federal government would rescue them, the states bound their own hands. American
state debt would be honored, when it was honored, by choice under market
discipline---a nineteenth-century anticipation of how sovereign debt without a
supranational enforcer is disciplined still.

\section{Relevance to the Book}
\label{sec:relevance}

Three threads seem worth carrying into the manuscript. First, the \emph{continuity
of the sovereign-immunity question}: \emph{Chisholm} arose from a Revolutionary
supply debt, the very category our companion note follows through the Treaty of
Paris and \emph{Ware v. Hylton}---but where a \emph{treaty} obligation could be
enforced against a state's law in federal court, a state's \emph{own} debt could
not, once the Eleventh Amendment intervened. Second, the \emph{contrast between
federal and state credit}: the United States built enforceable, reputation-backed
credit after 1790, while the states operated in a regime where only reputation
disciplined borrowing---and the 1840s show what that regime produced. Third, the
1840s as a \emph{clean test of reputational sovereign-debt theory}, with the
repay/repudiate split (Table~\ref{tab:defaults}) and the subsequent market
differentiation as the data. Each connects to material already drafted on the
Revolutionary debts, on Hamilton's assumption, and on the constitutional debt
reforms.

\section{Caveats and Verification Notes}
\label{sec:caveats}

\begin{itemize}
  \item \textbf{The Rothschild ``not a dollar'' quotation.} The wording is well
    attested and is best cited to \citet[39--40]{sexton2005debtor}, but the
    identity of the American agent addressed (often given as Duff Green) is
    unsettled in the secondary literature; confirm the name in Sexton or McGrane
    before printing it.
  \item \textbf{The 1842 federal loan.} That the federal government failed to place
    a loan abroad in 1842 is well attested; a specific dollar figure (sometimes
    given as \$12 million) we could not confirm and have omitted---check the loan
    act in the U.S. \emph{Statutes at Large}.
  \item \textbf{Sydney Smith.} The ``seize and divide'' passage is pinned to
    \citet[59]{mcgrane1935foreign} (reprinting the \emph{Times}, 4~Nov.\ 1843); the
    petition line to \citet[9]{sydneysmith1843letters}. The often-quoted aphorism
    ``there is no distinction between a nation that will not pay and an individual
    who will not pay'' we could not verify and do not attribute to him.
  \item \textbf{The foreign-held share of state debt.} Contemporary estimates
    ranged widely---from about \$65 million (1838) to \$110--165 million (1839),
    with British holdings usually put near \$100 million---so ``perhaps half or
    more held abroad'' is a characterization, not a single sourced figure.
  \item \textbf{The count of defaulters.} It was eight states plus the Florida
    Territory (nine governments). Ohio did \emph{not} default, though it is
    sometimes wrongly listed.
  \item \textbf{Monograph and article page numbers.} Specific pages in McGrane,
    Ratchford, Orth, and the verbatim wording of the English (1996) abstract were
    not all pinned to the physical volumes and should be checked before quotation.
\end{itemize}

\bibliographystyle{aer}
\bibliography{eleventh_amendment_state_debts}

\end{document}
